\chapter{Fundamentals of Graph Theory}
\label{ch:graph-theory}
\minitoc

The applications surveyed in Chapter~\ref{ch:applications} share a single
underlying object: the graph. To move from describing tasks to formulating
them precisely, that object must itself be defined formally. This
chapter introduces the formal language of graphs used throughout the rest of
this work. It first defines graphs, neighborhoods, and node and edge features;
it then encodes graph structure algebraically through the adjacency
matrix and distinguishes the main types of graphs; it examines the structural
properties, from walks to isomorphism, on which later developments rely; it
introduces the graph Laplacian and the spectral perspective it provides; and
it closes with a summary of the notation fixed here.

\section{Formal Definitions}
\label{sec:formal-definitions}

A \textit{graph} is a pair
\begin{equation}
  \mathcal{G} = (V, E),
  \label{eq:graph-definition}
\end{equation}
where $V$ is a finite set of \textit{nodes}, also called vertices, and $E$
is a set of \textit{edges}, also called links, each one an unordered pair of
nodes recording which two nodes are connected~\cite{diestel2017}. The number
of nodes is denoted $n = \lvert V \rvert$ and the number of edges
$m = \lvert E \rvert$. In classical graph-theoretic notation, generic nodes
are denoted by $u, v \in V$, and an edge between them by the unordered pair
$\{u, v\} \in E$; at this stage, the notation simply states that the two
nodes are connected.

Because $V$ and $E$ are sets in the mathematical sense, neither carries an
intrinsic order: a graph specifies which nodes are connected, and nothing
more. It is nevertheless convenient, particularly when graphs are
represented by matrices, to fix an indexing of the nodes by integers.
Throughout this work, nodes are therefore denoted
$i, j \in \{1, \ldots, n\}$, with the convention that $i$ refers to the node
under consideration and $j$ to a neighboring node whenever a neighborhood
relation is involved. The choice of indexing is arbitrary and carries no
information about the graph itself, a point this chapter returns to.

Two nodes are \textit{adjacent}, or \textit{neighbors}, when an edge
connects them. The \textit{neighborhood} of a node $i$ is the set of its
neighbors,
\begin{equation}
  \mathcal{N}(i) = \{\, j \in V : \{i, j\} \in E \,\},
  \label{eq:neighborhood}
\end{equation}
and the \textit{degree} of node $i$ is the number of neighbors it has,
\begin{equation}
  \deg(i) = \lvert \mathcal{N}(i) \rvert.
  \label{eq:degree}
\end{equation}
In a social network, for example, the neighborhood of a person is the set
of people connected to them, and the degree counts them.

Beyond its structure, a graph typically carries data~\cite{hamilton2020}.
Each node $i$ has an associated \textit{node feature vector}
$\mathbf{x}_i \in \mathbb{R}^d$, where $d$ is the node feature dimension,
describing the attributes of the entity the node represents, such as the
textual content of a paper in a citation network, or the chemical element of
an atom in a molecule. Stacking the $n$ node feature vectors as rows
produces the \textit{node feature matrix}
\begin{equation}
  \mathbf{X} =
  \begin{bmatrix}
    \mathbf{x}_1^\top \\
    \vdots \\
    \mathbf{x}_n^\top
  \end{bmatrix}
  \in \mathbb{R}^{n \times d},
  \label{eq:node-features}
\end{equation}
whose $i$-th row is the transposed feature vector of node $i$. Writing the
features as a matrix presupposes that the node indexing is fixed as above.
Edges can carry attributes as well: the \textit{edge feature vector}
$\mathbf{e}_{i,j} \in \mathbb{R}^c$, where $c$ is the edge feature
dimension, describes the connection between nodes $i$ and $j$, such as the
type of a chemical bond in a molecule.

\section{The Adjacency Matrix and Graph Types}
\label{sec:adjacency-types}

With the nodes indexed from $1$ to $n$, the structure of a graph can be
encoded in a single matrix. The \textit{adjacency matrix} of $\mathcal{G}$
is the matrix $\mathbf{A} \in \mathbb{R}^{n \times n}$ with entries
\begin{equation}
  A_{ij} =
  \begin{cases}
    1 & \text{if } \{i, j\} \in E, \\
    0 & \text{otherwise},
  \end{cases}
  \label{eq:adjacency-matrix}
\end{equation}
so that each entry records the presence or absence of a single edge.

The adjacency matrix depends on the node
indexing.\marginnote{There are $n!$ possible indexings of $n$ nodes, and
each yields a valid adjacency matrix of the same graph.} Reordering the
nodes permutes the rows and columns of $\mathbf{A}$, producing a different
matrix that describes exactly the same graph.
Figure~\ref{fig:adjacency-matrix} shows a graph together with its adjacency
matrices under two different indexings: a graph has not one adjacency matrix
but one per indexing, and no single one of them is privileged.

\begin{figure*}[t]
  \centering
  \includegraphics[width=\textwidth]{adjacency-matrix}
  \caption[A graph and its adjacency matrices under two node indexings.]{%
    A graph on six nodes (left) and its adjacency matrices under two
    different indexings of the nodes (center and right). Both matrices
    describe the same structure: reordering the nodes permutes the rows and
    columns of $\mathbf{A}$.
  }
  \label{fig:adjacency-matrix}
\end{figure*}

The first distinction among graph types concerns the orientation of edges.
In an \textit{undirected} graph, the setting assumed so far, an edge
$\{i, j\}$ has no orientation and can be traversed in either direction; its
adjacency matrix is consequently symmetric, $\mathbf{A} = \mathbf{A}^\top$.
In a \textit{directed} graph, edges are instead ordered pairs: the pair
$(i, j) \in E$ denotes an edge from a \textit{source} node $i$ to a
\textit{target} node $j$, and the presence of $(i, j)$ no longer implies the
presence of $(j, i)$. The adjacency entry reads $A_{ij} = 1$ exactly when
$(i, j) \in E$, so $\mathbf{A}$ is in general not symmetric and row $i$ lists
the targets of the edges leaving node $i$. The degree of a node accordingly
splits in two: its \textit{out-degree}, the number of edges leaving it, and
its \textit{in-degree}, the number of edges arriving at it. Unless stated
otherwise, the graphs treated in this work are undirected.

A second distinction concerns whether edges carry a magnitude. In a
\textit{weighted} graph, each edge has an associated scalar weight
$w_{ij} \in \mathbb{R}$, expressing, for instance, the strength of a
relationship or a physical distance; the adjacency matrix then stores
$A_{ij} = w_{ij}$ in place of $1$ (Figure~\ref{fig:directed-weighted}). In an
\textit{unweighted} graph, every edge counts the same and
\eqref{eq:adjacency-matrix} applies as stated. A scalar weight can be
viewed as the special case $c = 1$ of the edge feature vectors introduced
in Section~\ref{sec:formal-definitions}, that is, a single number attached
to each edge.

\begin{figure}[h]
  \centering
  \includegraphics[width=\textwidth]{directed-weighted}
  \caption[A directed graph and a weighted graph.]{%
    Left: a directed version of the graph of
    Figure~\ref{fig:adjacency-matrix} and its adjacency matrix. Edge
    orientations make $\mathbf{A}$ asymmetric: row $i$ lists the targets of
    the edges leaving node $i$, and the in-degree and out-degree of one
    highlighted node are annotated. Right: a weighted graph, with edge
    thickness proportional to the weight $w_{ij}$.
  }
  \label{fig:directed-weighted}
\end{figure}

A \textit{self-loop} is an edge connecting a node to itself; in the
adjacency matrix it appears as a nonzero diagonal entry, $A_{ii} \neq 0$. A
\textit{multigraph} allows several parallel edges between the same pair of
nodes, a multiplicity that the binary entries of \eqref{eq:adjacency-matrix}
cannot record. A graph with no self-loops and at most one edge between any
pair of nodes is called \textit{simple}. A \textit{bipartite} graph, in
contrast, is not a departure from simplicity but a structural subclass of
it: its nodes partition into two disjoint subsets, and every edge connects a node
of one subset with a node of the other, as in a network of users and products
joined by purchases. Because of this specificity, such graphs are explicitly 
named bipartite; a graph described only as simple implies no such partition.
Figure~\ref{fig:nonsimple-bipartite} illustrates the three cases.
As established in Chapter~\ref{ch:introduction}, the graphs treated
in this work are simple.

\begin{figure*}[t]
  \centering
  \includegraphics[width=\textwidth]{nonsimple-bipartite}
  \caption[Self-loops, a multigraph, and a bipartite graph.]{%
    Left: a graph with self-loops. Center: a multigraph with parallel
    edges. Right: a bipartite graph, with its two subsets of nodes drawn as
    distinct shapes.
  }
  \label{fig:nonsimple-bipartite}
\end{figure*}

The last distinction concerns the kinds of entities and relations a graph
contains.\marginnote{The homogeneous/heterogeneous terminology comes from
the network analysis and machine learning literature rather than from
classical graph theory.} A graph is \textit{homogeneous} when all of its
nodes represent one kind of entity and all of its edges represent one kind
of relation: every node then shares the same feature space $\mathbb{R}^d$
and every edge the same feature space $\mathbb{R}^c$. A citation network is
homogeneous, since every node is a paper and every edge a citation. A graph
is \textit{heterogeneous} when it contains several types of nodes or edges,
each with its own semantics and, typically, its own feature
space~\cite{hamilton2020}. A citation network that includes both papers and
their authors is heterogeneous: papers and authors are different kinds of
entities, joined by a relation of authorship.
Figure~\ref{fig:homogeneous-heterogeneous} contrasts the two settings. As
stated in Chapter~\ref{ch:introduction}, this work treats homogeneous
graphs.

\begin{figure}
  \centering
  \includegraphics[width=\textwidth]{homogeneous-heterogeneous}
  \caption[A homogeneous and a heterogeneous graph.]{%
    Left: a homogeneous graph, in which all nodes represent the same kind
    of entity and all edges the same relation. Right: a heterogeneous
    graph, in which node shapes denote different entity types and edge
    styles denote different relations.
  }
  \label{fig:homogeneous-heterogeneous}
\end{figure}

\section{Walks, Connectivity, Homophily, and Isomorphism}
\label{sec:structural-properties}

The definitions so far describe a graph locally: nodes, their neighbors,
and their attributes. Several properties that matter in the remainder of
this work are instead global, arising from how edges chain together across
the graph. A \textit{walk} is a sequence of nodes in which each consecutive
pair is joined by an edge, and its \textit{length} is the number of edges
traversed; a walk may revisit nodes and edges. A \textit{path} is a walk in
which all nodes are distinct (Figure~\ref{fig:walks-paths}).

\begin{marginfigure}
  \includegraphics{walks-paths}
  \caption[A walk and a shortest path.]{%
    A walk that revisits a node (top) and a shortest path between the same
    endpoints (bottom), whose length is their geodesic distance.
  }
  \label{fig:walks-paths}
\end{marginfigure}

The \textit{geodesic distance} between two nodes is the length of a
shortest path between them,
\begin{equation}
\begin{aligned}
d(i,j)
=
\min \Bigl\{
  k \in \mathbb{N} : {}&
  \text{there is a path of length } k \\
  &\text{between } i \text{ and } j
\Bigr\}
\end{aligned}
\label{eq:geodesic-distance}
\end{equation}
with $d(i, i) = 0$. The distance is always written with its two arguments,
$d(i, j)$, which distinguishes it from the node feature dimension $d$. When
no path joins $i$ and $j$, the distance is taken to be infinite, a
situation characterized next.

A graph is \textit{connected} when a path exists between every pair of its
nodes, that is, when $d(i, j)$ is finite for all $i, j \in V$. When it is
not, the graph decomposes into \textit{connected components}: maximal
groups of nodes that are mutually reachable, with no edges between groups
(Figure~\ref{fig:connectivity}). In directed graphs, connectivity refines
into weak and strong variants according to whether edge orientations are
respected; since this work treats undirected graphs, the refinement is not
developed further.

\begin{figure}
  \centering
  \includegraphics[width=\textwidth]{connectivity}
  \caption[A disconnected graph and its components.]{%
    A disconnected graph with two connected components, one of which is a
    single isolated node.
  }
  \label{fig:connectivity}
\end{figure}

Structure and attributes are often correlated. A graph exhibits
\textit{homophily} when adjacent nodes tend to have similar attributes,
and \textit{heterophily} when they tend to differ. The notion
originates in the analysis of social networks, where similar individuals
form ties at higher rates than dissimilar ones~\cite{mcpherson2001}, and it
extends to many of the graphs of Chapter~\ref{ch:applications}: in a
citation network, a paper predominantly cites papers on related topics, so
neighboring nodes tend to share a subject. Whether a graph is homophilous
or heterophilous is an empirical property of the data rather than a
structural constraint, and the degree to which it holds varies from one
graph to another.

The arbitrariness of node indexing observed in
Section~\ref{sec:adjacency-types} raises a final question: when do two
graphs that are presented differently have the same structure? Two graphs
$\mathcal{G}_1 = (V_1, E_1)$ and $\mathcal{G}_2 = (V_2, E_2)$ are
\textit{isomorphic} when there exists a bijection $\pi : V_1 \to V_2$ that
preserves adjacency,
\begin{equation}
 \{i, j\} \in E_1
  \quad \Longleftrightarrow \quad
  \{\pi(i), \pi(j)\} \in E_2,
  \label{eq:graph-isomorphism}
\end{equation}
for all nodes $i, j \in V_1$. Isomorphic graphs are the same graph up to a
relabeling of its nodes; the two adjacency matrices of
Figure~\ref{fig:adjacency-matrix} are representations related in exactly
this way. Two graphs drawn very differently may be isomorphic
(Figure~\ref{fig:isomorphic-graphs}). Determining whether two graphs
are isomorphic is a fundamental problem in graph theory.
A classical heuristic, the Weisfeiler--Leman test, distinguishes
most non-isomorphic graphs in practice~\cite{morris2023}, and it reappears later in this work
in connection with the expressive power of graph neural networks.

\begin{marginfigure}
  \includegraphics{isomorphic-graphs}
  \caption[Two isomorphic graphs.]{%
    Two graphs that are drawn differently but are isomorphic: a bijection
    between their node sets preserves all adjacencies.
  }
  \label{fig:isomorphic-graphs}
\end{marginfigure}

\section{The Graph Laplacian and the Spectral Perspective}
\label{sec:laplacian-spectral}

The structural properties of Section~\ref{sec:structural-properties} and the
algebraic encoding of Section~\ref{sec:adjacency-types} meet in a single
observation: powers of the adjacency matrix count walks. For any positive
integer $k$, the entry in row $i$ and column $j$ of the $k$-th matrix power
gives the number of walks of length $k$ between the two nodes,
\begin{equation}
  (\mathbf{A}^k)_{ij} =
  \bigl\lvert \{\, \text{walks of length } k \text{ between } i
  \text{ and } j \,\} \bigr\rvert,
  \label{eq:walk-counting}
\end{equation}
a count that is symmetric in $i$ and $j$ for the undirected graphs treated
here. In particular, the entry is nonzero exactly when the two nodes are
joined by some walk of that length, which ties the reachability notions of
Section~\ref{sec:structural-properties} directly to matrix algebra and
motivates the matrices introduced next.

The \textit{degree matrix} $\mathbf{D} \in \mathbb{R}^{n \times n}$ is the
diagonal matrix that holds the node degrees,
\begin{equation}
  \mathbf{D} = \operatorname{diag}\bigl(\deg(1), \ldots, \deg(n)\bigr),
  \label{eq:degree-matrix}
\end{equation}
with every off-diagonal entry equal to zero. Combining it with the adjacency
matrix yields the central operator of spectral graph
theory~\cite{chung1997}, the \textit{graph Laplacian}
\begin{equation}
  \mathbf{L} = \mathbf{D} - \mathbf{A}.
  \label{eq:laplacian}
\end{equation}
Each diagonal entry of $\mathbf{L}$ is the degree of the corresponding node,
and each off-diagonal entry is $-1$ where an edge is present and $0$
otherwise. For the six-node graph of Figure~\ref{fig:adjacency-matrix},
under the node indexing of its centre panel and with
$\mathbf{D} = \operatorname{diag}(2, 3, 2, 3, 3, 1)$, the Laplacian is
\[
  \mathbf{L} =
  \begin{bmatrix}
    \phantom{-}2 & -1 & \phantom{-}0 & \phantom{-}0 & -1 & \phantom{-}0 \\
    -1 & \phantom{-}3 & -1 & \phantom{-}0 & -1 & \phantom{-}0 \\
    \phantom{-}0 & -1 & \phantom{-}2 & -1 & \phantom{-}0 & \phantom{-}0 \\
    \phantom{-}0 & \phantom{-}0 & -1 & \phantom{-}3 & -1 & -1 \\
    -1 & -1 & \phantom{-}0 & -1 & \phantom{-}3 & \phantom{-}0 \\
    \phantom{-}0 & \phantom{-}0 & \phantom{-}0 & -1 & \phantom{-}0 & \phantom{-}1
  \end{bmatrix}.
\]

The Laplacian admits a direct interpretation in terms of how values spread
over a graph. Assigning a real number to each node defines a vector
$\mathbf{f} \in \mathbb{R}^n$ called a \textit{graph signal}, 
with component $f_i$ the value at node $i$. The Laplacian measures how much
such an assignment varies across edges through the quadratic form
\begin{equation}
  \mathbf{f}^\top \mathbf{L} \mathbf{f}
  = \sum_{\{i, j\} \in E} (f_i - f_j)^2,
  \label{eq:laplacian-quadratic-form}
\end{equation}
which sums the squared difference between the values at the two endpoints of
every edge. The form vanishes when $\mathbf{f}$ is constant on each connected
component and grows as the values assigned to neighboring nodes diverge.
Because it is a sum of squares, the quadratic form is nonnegative for every
$\mathbf{f}$, so the Laplacian is \textit{positive semidefinite}.

Being real and symmetric, the Laplacian admits an eigendecomposition
\begin{equation}
  \mathbf{L} = \mathbf{U} \boldsymbol{\Lambda} \mathbf{U}^\top,
  \label{eq:eigendecomposition}
\end{equation}
where the columns of $\mathbf{U} \in \mathbb{R}^{n \times n}$ form an
orthonormal basis of eigenvectors and
$\boldsymbol{\Lambda} = \operatorname{diag}(\lambda_1, \ldots, \lambda_n)$
collects the corresponding eigenvalues. Symmetry guarantees that these
eigenvalues are real, and positive semidefiniteness that they are
nonnegative; ordered, they satisfy
$0 = \lambda_1 \leq \cdots \leq \lambda_n$. The smallest eigenvalue is always
zero, with the constant vector as an eigenvector, since the rows of
$\mathbf{L}$ sum to zero. More can be said: the multiplicity of the
eigenvalue $0$ equals the number of connected components of the graph,
linking the spectral description back to the connectivity of
Section~\ref{sec:structural-properties}.

For many purposes the Laplacian is rescaled by the node degrees. The
\textit{symmetric normalized Laplacian} is
\begin{equation}
  \mathbf{L}_{\mathrm{sym}}
  = \mathbf{D}^{-1/2} \mathbf{L} \mathbf{D}^{-1/2}
  = \mathbf{I} - \mathbf{D}^{-1/2} \mathbf{A} \mathbf{D}^{-1/2},
  \label{eq:normalized-laplacian}
\end{equation}
where $\mathbf{I}$ is the identity matrix and $\mathbf{D}^{-1/2}$ is the
diagonal matrix of inverse square-root
degrees.\marginnote{For a node of degree zero the corresponding entry of
$\mathbf{D}^{-1/2}$ is taken to be zero by convention.} It is again symmetric
and positive semidefinite, so it shares the spectral structure just
described, while its eigenvalues are confined to the interval
$[0, 2]$~\cite{chung1997}. This normalized form is the variant used in the
spectral developments of later chapters. Together, these properties make the
Laplacian a natural bridge between the structure of a graph and the linear
algebra used to operate on it.

\FloatBarrier
\section{Notation Summary}
\label{sec:notation-summary}

The notation introduced in this chapter is collected in
Table~\ref{tab:notation} for reference. It is used consistently throughout
the remainder of this work.

\begin{table}
  \centering
  \caption[Notation introduced in this chapter.]{%
    Notation introduced in this chapter, with the equation or section in
    which each symbol is defined.
  }
  \label{tab:notation}
  \begin{tabular}{@{}lll@{}}
    \toprule
    Symbol & Meaning & Reference \\
    \midrule
    $\mathcal{G} = (V, E)$ & Graph with node set $V$ and edge set $E$
      & \eqref{eq:graph-definition} \\
    $n,\ m$ & Number of nodes and edges
      & \S\ref{sec:formal-definitions} \\
    $i,\ j$ & Nodes; $i$ central, $j$ a neighbor
      & \S\ref{sec:formal-definitions} \\
    $\mathcal{N}(i)$ & Neighborhood of node $i$
      & \eqref{eq:neighborhood} \\
    $\deg(i)$ & Degree of node $i$
      & \eqref{eq:degree} \\
    $d(i, j)$ & Geodesic distance between $i$ and $j$
      & \eqref{eq:geodesic-distance} \\
    $\pi$ & Isomorphism bijection between node sets
      & \eqref{eq:graph-isomorphism} \\
    \midrule
    $\mathbf{x}_i,\ d$ & Node feature vector and its dimension
      & \eqref{eq:node-features} \\
    $\mathbf{X}$ & Node feature matrix
      & \eqref{eq:node-features} \\
    $\mathbf{e}_{i,j},\ c$ & Edge feature vector and its dimension
      & \S\ref{sec:formal-definitions} \\
    $w_{ij}$ & Scalar edge weight
      & \S\ref{sec:adjacency-types} \\
    \midrule
    $\mathbf{A}$ & Adjacency matrix
      & \eqref{eq:adjacency-matrix} \\
    $(\mathbf{A}^k)_{ij}$ & Walks of length $k$ between $i$ and $j$
      & \eqref{eq:walk-counting} \\
    $\mathbf{D}$ & Degree matrix
      & \eqref{eq:degree-matrix} \\
    $\mathbf{L}$ & Graph Laplacian
      & \eqref{eq:laplacian} \\
    $\mathbf{L}_{\mathrm{sym}}$ & Symmetric normalized Laplacian
      & \eqref{eq:normalized-laplacian} \\
    \midrule
    $\mathbf{f}$ & Graph signal, one real value per node
      & \eqref{eq:laplacian-quadratic-form} \\
    $\mathbf{U},\ \boldsymbol{\Lambda}$ & Eigenvectors and eigenvalues of
      $\mathbf{L}$ & \eqref{eq:eigendecomposition} \\
    $\lambda_1, \ldots, \lambda_n$ & Laplacian eigenvalues, in increasing
      order & \eqref{eq:eigendecomposition} \\
    \bottomrule
  \end{tabular}
\end{table}
